\section{Introduction}
\label{sec:Introduction}

Archimesh is a collaborative real-time modelling server for Archi, the ArchiMate
modelling tool. It lets multiple modellers work on the same model simultaneously,
with the server acting as the authoritative source of truth for the shared model
timeline.

This manual is written for developers and operators who need to set up, run,
supervise and troubleshoot an Archimesh deployment. It assumes familiarity with
Java, Docker, Neo4j and Kafka at a working level.

\subsection{What Archimesh offers}

Archimesh provides five core capabilities:

\begin{itemize}
\item \textbf{Real-time collaborative editing.} Multiple Archi clients connect to
  one server, capture local edits as structured operations, and apply remote
  operations as they are committed. The server assigns a single authoritative
  revision order so every client converges on the same model state.

\item \textbf{Authoritative model storage.} Neo4j stores both the current
  materialized model graph and the append-only operation log. The server is the
  source of truth; local Archi files are treated as cache projections.

\item \textbf{CRDT-informed merge.} The server applies Last-Writer-Wins (LWW)
  per-field merges with causal clocks and tombstone protection for concurrent
  edits on the same diagram objects. Lease-based locks protect noisy notation
  edits (resize, move, bendpoints, style) from destructive interleaving.

\item \textbf{Offline resilience.} When a client loses its connection, local
  edits are queued in a bounded persistent outbox and replayed in FIFO order
  after reconnect. Replayed ops are rebased to the latest server revision before
  send, and conflicting stale entries are deterministically dropped.

\item \textbf{Operational visibility.} The admin API and admin UI expose model
  catalog, active sessions, recent activity, ACL management, audit trails,
  integrity checks, compaction control, and export/import for backup and
  migration.

\end{itemize}

\subsection{What this manual covers}

\begin{itemize}
\item Section~\ref{sec:GettingStarted} walks through building and running the
  server and client locally.
\item Section~\ref{sec:Architecture} describes the system architecture, component
  layout and major data flows.
\item Section~\ref{sec:ClientPlugin} covers the Archi client plugin: EMF change
  capture, operation mapping, remote apply, offline outbox and cache policy.
\item Section~\ref{sec:ServerApi} documents the WebSocket and REST API
  endpoints, message types and request/response contracts.
\item Section~\ref{sec:Persistence} explains the Neo4j graph model:
  materialized state, op-log, tags, tombstones and property clocks.
\item Section~\ref{sec:Messaging} covers the Kafka topic layout, partitioning
  requirements and fan-out path.
\item Section~\ref{sec:CrdtAndMerge} details the CRDT merge semantics:
  LWW field merges, tombstone protection, OR-set collections and lock model.
\item Section~\ref{sec:Authorization} documents identity modes, role
  normalization, per-model ACLs and the policy decision point.
\item Section~\ref{sec:AdminOperations} explains the admin UI and operational
  supervision: overview, models, versions, sessions, access and audit.
\item Section~\ref{sec:Management} covers system management: creating models,
  compactions, exports, imports, rebuilds, tags and integrity checks.
\item Section~\ref{sec:Troubleshooting} describes common failure modes, log
  analysis, diagnostic endpoints and recovery procedures.
\item Appendix~\ref{sec:AppendixScripts} is a reference for all scripts in the
  \filepath{scripts/} directory.
\item Appendix~\ref{sec:AppendixSchemas} documents the wire protocol message
  types and JSON schema contracts.
\end{itemize}

\subsection{Terminology}

\begin{description}
\item[Model] A named linear revision timeline. One \texttt{HEAD}, zero or more
  immutable tags, no branches. Identified by a stable \texttt{modelId}.
\item[Operation (op)] A single Archimesh change: create, update or delete of an
  element, relationship, view, view object, connection, folder or property.
\item[Op batch] A client-submitted group of operations sharing one
  \texttt{opBatchId} for idempotency. The server assigns one contiguous
  revision range to the batch.
\item[Revision] A monotonically increasing timeline position. Each op in a
  batch consumes one revision. The revision range is the authoritative ordering.
\item[HEAD] The current tip of the model timeline. Only \texttt{HEAD} is
  writable; tagged revisions are read-only.
\item[Tag] An immutable named pointer to a historical revision, stored with a
  snapshot of the model state at that point.
\item[Materialized state] The current model graph in Neo4j, kept in sync with
  the op-log. Used for fast joins, exports and admin inspection.
\item[Outbox] A bounded persistent FIFO queue on the client side that holds
  local edits while disconnected from the server.
\item[Compaction] Admin-controlled pruning of old op-log commits and merge
  metadata below a configurable retention watermark. Tombstones are retained
  for delete safety.
\end{description}
